\begin{frame}{역사: 카지노 도시에 이름을 건 알고리즘}
  \begin{columns}[T]
    \begin{column}{0.55\textwidth}
      \begin{itemize}
        \item 이름의 기원: 모나코 공국의 카지노 도시
          \kb{몬테카를로}.
        \item 1940년대, MIT·맨해튼 프로젝트: Ulam,
          von Neumann, Metropolis (Ulam \& Metropolis, 1949).
        \item 풀어야 할 문제: \kb{중성자 이동} --- 물질 안에서
          중성자가 무작위로 여기저기 튕기는 과정.
        \item 이산 확률방정식을 정확히 풀어 계산하는 것은
          불가능에 가까움.
      \end{itemize}
    \end{column}
    \begin{column}{0.45\textwidth}
      \begin{block}{해법}
        중성자가 어디로 가는지 \textbf{주사위처럼 던져서}
        충분히 많이 추적해보자 --- \textbf{무작위성을 도구}로
        쓰는 계산법, \kb{몬테카를로 방법}.
      \end{block}
      \vskip0.6em
      \begin{block}{RL에서의 MC도 본질은 같음}
        전이확률을 몰라도, \textbf{무작위 전이를 실제로
        겪어본 기록(에피소드)의 평균}이 기댓값을 대신한다.
        카지노에서 주사위를 돌리는 것과 시뮬레이터에서
        에피소드를 돌리는 것은 \textbf{같은 행위}.
      \end{block}
    \end{column}
  \end{columns}
\end{frame}
